La contratación pública es un proceso central del gobierno electrónico: transforma los presupuestos públicos en contratos, obras y servicios mediante plataformas cada vez más digitales. También es relevante para la democracia electrónica, porque los registros abiertos de contratación amplían la capacidad de los organismos de control, periodistas, empresas y ciudadanos para examinar cómo se asignan los recursos públicos. Estas capacidades se alinean con las agendas de gobierno digital sobre datos abiertos y análisis responsables para la supervisión del sector público. Trabajos recientes en este sector presentan el análisis de contrataciones como una herramienta para la supervisión basada en riesgos \cite{salazar2024vigia}. Los documentos de licitación son fundamentales en este contexto porque definen las condiciones administrativas, los requisitos técnicos, las reglas de calificación y las restricciones de ejecución antes de que se conozca el resultado final. Evidencia reciente del ámbito de las contrataciones también muestra que las descripciones de las convocatorias permiten predecir resultados relacionados con la competencia y la eficiencia de costos \cite{acikalin2024procurementcalls}. Por tanto, la pregunta relevante no es solo si los resultados de las licitaciones pueden predecirse a partir del texto documental, sino también qué representación de ese texto aporta más información para un monitoreo responsable de la contratación pública.

Este trabajo estudia dicha pregunta mediante una comparación directa entre características léxicas dispersas y representaciones densas basadas en LLM. El análisis se concentra en licitaciones de construcción vial, cuyos documentos suelen ser relativamente estructurados y estandarizados. Los modelos léxicos dispersos continúan siendo líneas base sólidas e interpretables para la clasificación documental \cite{wang2012baselines}, y trabajos recientes sobre documentos extensos confirman que las representaciones léxicas o híbridas pueden seguir siendo competitivas, en lugar de quedar uniformemente superadas por arquitecturas transformer \cite{alvaprincipe2025survey}. Al mismo tiempo, la investigación sobre embeddings se ha extendido hacia representaciones densas basadas en LLM \cite{wang2024improvingembeddings}. En este contexto, el objetivo no es establecer una familia de modelos universalmente superior, sino comprender cuánta señal predictiva puede recuperarse de indicios léxicos, cuánta ya está capturada por embeddings densos congelados y si el modelado explícito de interacciones entre fragmentos aporta valor adicional.

La canalización utiliza un LLM causal preentrenado como codificador congelado para producir embeddings por fragmento a partir de documentos de licitación en español. Se comparan dos estrategias simples de agregación de estos embeddings, basadas en agrupación por promedio, con un transformer liviano que modela las interacciones entre fragmentos. Estos enfoques basados en LLM se evalúan junto con TF--IDF más regresión logística y líneas base triviales. Este diseño permite separar los efectos de la representación de los efectos del modelado posterior.

El estudio ofrece una comparación reproducible entre líneas base léxicas, triviales, de embeddings agregados y de un transformer entre fragmentos, bajo validación cruzada determinista de cinco pliegues. También analiza por separado el desempeño de ordenamiento, la calibración y la sensibilidad al umbral. El resultado central presenta matices: TF--IDF más regresión logística es el modelo con mejor ordenamiento, mientras que el transformer entre fragmentos ofrece la mejor calibración y el comportamiento más estable ante cambios del umbral. Esto aporta evidencia útil para la pregunta más amplia de cómo se distribuye la señal textual en los documentos de contratación pública. Los artefactos asociados están archivados en Zenodo (doi:10.5281/zenodo.20128769).
